% Generated draft: requires mathematical and source review.
\documentclass[11pt]{article}
\usepackage[margin=25mm]{geometry}
\usepackage{fontspec}
\setmainfont{Latin Modern Roman}
\usepackage{amsmath,amssymb,mathtools,unicode-math}
\setmathfont{Latin Modern Math}
\usepackage{enumitem,needspace,xurl,hyperref}
\setlength{\parindent}{0pt}
\setlength{\parskip}{5pt}
\setlength{\emergencystretch}{3em}
\newcommand{\N}{\mathbb{N}}
\newcommand{\Z}{\mathbb{Z}}
\newcommand{\Q}{\mathbb{Q}}
\newcommand{\R}{\mathbb{R}}
\newcommand{\C}{\mathbb{C}}
\newcommand{\F}{\mathbb{F}}
\newcommand{\A}{\mathbb{A}}
\newcommand{\PP}{\mathbb{P}}
\newcommand{\E}{\mathbb{E}}
\newcommand{\eps}{\varepsilon}
\newcommand{\dd}{\,\mathrm{d}}
\begin{document}
\section*{UnsolvedMath ID: 30001369}
\textbf{Source problem code:} catalogue entry
\subsection{Definitions and mathematical statement}
The supplied problem concerns piecewise-affine dynamical systems preserving complementary stable and unstable cone fields. Write \(T\colon X\to X\) for such a system. Piecewise affine means that, on each element of a specified partition of the phase space, the map is affine in the relevant affine coordinates. The record does not specify the phase space, partition, treatment of partition boundaries, or the precise invariance, complementarity, and expansion or contraction conditions for the cone fields. Consequently, it does not determine the class of systems being quantified over.

For reference, under the unweighted discrete-time periodic-point convention, put
\[
\operatorname{Fix}(T^n)=\{x\in X:T^n x=x\},\qquad N_n=\#\operatorname{Fix}(T^n),
\]
and define
\[
\zeta_T(z)=\exp\!\left(\sum_{n=1}^{\infty}\frac{N_n}{n}z^n\right).
\]
This formula defines a holomorphic germ at \(0\) if all \(N_n\) are finite and the displayed series has positive radius of convergence. These conditions and this particular zeta-function convention are not specified in the supplied record; they are conditional explanations of the notation, not additional source hypotheses.

\textbf{Requested statement, subject to the unresolved specifications.} For each piecewise-affine system preserving complementary stable and unstable cone fields in the class intended by the source report, construct a meromorphic continuation of its Ruelle zeta function and determine domains on which that continuation exists. More precisely, once its defining holomorphic germ on a neighborhood \(V\) of \(0\) is fixed, seek a connected open domain \(U\subseteq\C\) containing \(V\) and a meromorphic function \(Z_T\) on \(U\) satisfying
\[
Z_T|_V=\zeta_T|_V.
\]
The intended continuation should extend beyond the original domain of definition. The supplied record does not specify a target domain, a quantitative lower bound for continuation, or whether maximal domains are sought. Nor does it identify the exact Ruelle zeta convention. A fully quantified formulation therefore requires the missing definitions from the report; the displayed periodic-point formula cannot be asserted to be its intended normalization.
\subsection{Short English statement}
For the piecewise-affine systems with stable and unstable cone fields intended by the report, construct meromorphic continuations of their Ruelle zeta functions and determine where they continue. The supplied record leaves the precise system class, zeta convention, and target domains unresolved. Its \texttt{partially\_solved} label is supplied dated curation, not an independently verified conclusion.
\subsection{Sources}
\begin{itemize}
\item Supplied local project registry record, permanent UnsolvedMath ID \(30001369\), source problem code \texttt{catalogue entry}, title: \emph{Meromorphic Ruelle Zeta Functions for Piecewise Hyperbolic Systems}. The record belongs to the Oberwolfach Reports Open Problems collection. No individual proposer is supplied.
\item Supplied original report citation: \emph{Mini-Workshop: Spectrum of Transfer Operators -- Recent Developments and Applications} (2010), Oberwolfach Reports, pp.~2699--2730, DOI \texttt{10.4171/owr/2009/49}. \url{https://doi.org/10.4171/owr/2009/49}. The literature-source label is ``original Oberwolfach report 4132.'' The record separately gives proposed year \(2009\). Its original question asks about meromorphic extensions for the ``piecewise hyperbolic systems under discussion.''
\item Supplied literature source labeled ``weighted zeta continuation theorem'': \url{https://doi.org/10.1007/s00220-022-04538-z}.
\item Supplied literature source labeled ``Ruelle zeta survey'': \url{https://smf.emath.fr/system/files/filepdf/AST_1990__187-188__1_0.pdf}.
\item Literature assessment dated exactly \texttt{2026-08-21}: the supplied curation labels the problem partially solved and states that meromorphic continuation is known in several smooth, expanding, and open-hyperbolic settings, while the unrestricted piecewise-affine systems in the report are not covered by a general theorem. This is supplied dated curation, not independently verified here; its time-sensitive status and cited claims require independent verification.
\item The registry supplies statement status \texttt{corrected\_verified} and says that the clean statement was checked against the official report and preserves its hypotheses, notation, and open target. This verification claim is source metadata and has not been independently verified here.
\end{itemize}
\end{document}
